\chapter*{ЗАКЛЮЧЕНИЕ}
\addcontentsline{toc}{chapter}{ЗАКЛЮЧЕНИЕ}

В ходе выполнения курсовой работы была успешно достигнута поставленная цель: разработан программный комплекс для трёхмерного моделирования и визуализации движения автомобиля в условиях городской застройки.

В процессе решения поставленных задач были получены результаты.

\begin{enumerate}
	\item \textbf{Проведен анализ} предметной области и алгоритмической базы компьютерной графики. Обоснован выбор алгоритма построчного сканирования с использованием Z-буфера как наиболее надежного метода удаления невидимых поверхностей для сцен произвольной сложности.
	
	\item \textbf{Спроектирована и реализована архитектура} графического приложения, включающая модули управления сценой, математического ядра и подсистему рендеринга. Выбранный стек технологий (Python + NumPy) позволил реализовать векторизованные вычисления на центральном процессоре.
	
	\item \textbf{Реализован графический конвейер}. В приложение внедрены следующие методы визуализации:
	\begin{itemize}
		\item построение динамических теней методом карт теней, что обеспечивает корректное самозатенение объектов;
		\item плоское закрашивание с использованием модели освещения Ламберта;
		\item UV-развертка для наложения текстур для повышения визуальной детализации дорожного полотна и зданий.
	\end{itemize}
	
	\item \textbf{Разработан модуль навигации.} Реализован модифицированный волновой алгоритм поиска пути на ориентированном графе, который учитывает не только геометрию препятствий, но и правила дорожного движения (разрешенные направления поворотов), задаваемые специальной матрицей.
	
	\item \textbf{Выполнено исследование} производительности системы. Установлено, что фактором, влияющим на время генерации кадра, является коэффициент масштабирования сцены.
\end{enumerate}

Разработанное программное обеспечение представляет собой законченную систему, демонстрирующую работу фундаментальных алгоритмов компьютерной графики и поиска пути. Полученные результаты могут быть использованы в учебных целях для демонстрации принципов работы графических движков, а также как база для создания симуляторов транспортных потоков.